\section{What one affine optimum does not encode}
\label{sec:boundaries}

\subsection{Distinct framed states still exist}

The previous regular-cover gluing report~\cite{certificates} counts
$m^\beta$ possible cyclic cycle-phase tuples after tree normalization.
That count concerns distinct coverings over a fixed labelled
decomposition. It cannot be removed when the requested output is an
explicit list of every such state. Moreover, its local block covers are
already connected, so a connected assembly of those blocks is connected
for every seam phase. In that exact model, connectivity optimization
would be trivial.

The current theorem concerns a different query on a broader input model:
possibly disconnected local cyclic covers, with affine dependence in
their monodromy and gluing data. It computes the optimum of one invariant
over a family. It neither identifies all members with one another nor
provides a complete invariant for all future attachments. This distinction
prevents an incorrect inference from a polynomial existential query to
a polynomial enumeration or canonicalization of all geometric states.

\subsection{Connectedness does not determine genus}

\begin{example}[A five-sheet pair of pants]
\label{ex:genus}
Let the base be a pair of pants, and let its first two boundary
translations on $\Z/5\Z$ be $1$ and $a$. The third is $-1-a$.
For $a=1$ the cover is connected and has three boundary components,
one over each base boundary. Since its Euler characteristic is $-5$,
its genus is two. For $a=-1$ it remains connected but has seven boundary
components: one, one, and five. Its genus is zero.
\end{example}

The example follows directly from orbit lengths of a translation and
$\chi(\widetilde F)=5\chi(F)$. It is also checked by the existing cover
classifier in the regression suite. A hierarchy step that needs a planar
base, a disc, a specific peripheral curve, or a particular attachment
type therefore cannot substitute the component optimizer for that
geometric test.

\subsection{Separate component optima can conflict}

Even a sum of simple component-count objectives can carry correlations.
Over $\Z/2\Z$ consider two disconnected base graphs with respective
single loop voltages $z$ and $z+1$. Each graph individually admits a
connected cover. No common parameter makes both connected: their total
component counts are $1+2$ or $2+1$. Thus optimizing each component
separately and adding the answers would give the wrong result.
The current surface-family compiler requires a connected base graph.
A future multi-component interface must retain shared constraints and
prove its joint optimization rule.

\subsection{Counting can be harder than producing an optimum}

For this subsection only, assume the prime divisors of $m$ are known.
Let $S$ be the feasible parameter set, let $z_0+B u$ be a full
parameterization, and write $h=v+G u$, where $v=c+Hz_0$ and $G=HB$.
For each prime $p\mid m$, let $r_p$ be the rank of $G\pmod p$ and put
\[
 \theta_p=
 \begin{cases}
  1,&-v\notin\im(G\pmod p),\\
  1-p^{-r_p},&-v\in\im(G\pmod p).
 \end{cases}
\]

\begin{proposition}[Exact proportion of connected assignments]
\label{prop:counting}
The number of distinct feasible parameters giving a connected cyclic
cover is
\begin{equation}
 |S|\prod_{p\mid m}\theta_p.\label{eq:connected-count}
\end{equation}
\end{proposition}
\begin{proof}
Uniform $u\in(\Z/m\Z)^k$ pushes forward to uniform $z\in S$, because
every fibre of a homomorphism onto $S-z_0$ has the same cardinality.
Modulo a prime $p$, the affine equation $v+Gu=0$ has probability zero
when inconsistent and probability $p^{-r_p}$ otherwise. A vector is
primitive modulo $m$ exactly when it is not the zero vector modulo any
prime dividing $m$. The Chinese remainder theorem makes these prime
conditions independent for uniform $u$. Multiplying their probabilities
and then by $|S|$ proves the formula. Redundant parameterizations do not
alter the result.
\end{proof}

This is a counting identity, not a factor-free algorithm for evaluating
the product. The distinction is necessary. In the one-variable family
$h(z)=z$ with no constraints, the connected assignments are precisely
the units, so their number is Euler's totient $\varphi(m)$. If
$m=pq$ for distinct unknown primes, knowledge of $\varphi(m)$ gives
\[
 p+q=m-\varphi(m)+1,
\]
and hence recovers $p,q$ as the integer roots of a quadratic. Thus a
general polynomial-time exact connected-assignment counter would imply
a polynomial-time algorithm for factoring such semiprimes. This is a
reduction, not a proof that the counting problem has no polynomial
algorithm. It explains why the minimum and its witness can be obtained
without offering an equally strong counting claim.

\subsection{Many simultaneous peripheral conditions are hard}

\begin{theorem}[Simultaneous peripheral connectedness is NP-complete]
\label{thm:peripheral-hardness}
Consider the following decision problem: a connected planar bordered base
surface is supplied with a three-sheet cyclic cover whose generator
translations are affine functions of parameters in $\Z/3\Z$, and a
subset of boundary circles is specified. Does some parameter assignment
make the entire preimage of each specified boundary circle connected?
This problem is NP-complete, even with no additional congruence constraints.
\end{theorem}
\begin{proof}
Membership in NP follows by supplying the parameter vector. A boundary
translation $a$ has a connected three-sheet preimage exactly when
$a\ne0\pmod3$, so verification is polynomial in the listed input.

For hardness, take a finite graph $G=(V,E)$ and one parameter $z_v$ for
each vertex. Use a planar surface with $|E|+1$ boundary circles and
assign the first $|E|$ free generators the translations $z_u-z_v$,
one for each edge $uv$. The final boundary has the inverse product,
so its translation is minus the sum of these edge differences. Specify
only the first $|E|$ boundaries. Their preimages are all connected exactly
when $z_u\ne z_v$ for every edge, which is a proper three-colouring.
Graph three-colourability is NP-complete~\cite{gjs}. The construction is
polynomial even with dense coefficient rows. Restricting to graphs with
at least one edge preserves hardness, and in every satisfying instance
the entire surface cover is also connected, since at least one nonzero
translation generates $\Z/3\Z$.
\end{proof}

This is a hardness result for an expressive family query. It is not a
hardness result for unknot recognition, normal-surface embedding, or a
single already specified cover. It makes a practical interface distinction:
one global gcd objective is tractable, whereas imposing many correlated
peripheral objectives permits graph colouring. Consequently a future
hierarchy adapter must exploit geometric restrictions on its actual
families rather than assume arbitrary affine peripheral constraints are
equally inexpensive. The reduction also carries any quasi-polynomial
algorithm for this unrestricted peripheral query to graph three-colouring.
NP-completeness alone does not unconditionally rule out that running time;
it identifies the strength of the additional algorithm being requested.

\subsection{Completeness of a representation is a separate obligation}

For a specified normal surface, the extracted interval maps are usually
partial maps on several stacks of different sizes. They are not
automatically translations of a single cyclic fibre. Likewise a normal
surface search includes nonnegative coordinates and quadrilateral
alternatives, not merely unrestricted residue parameters. Passing from
those data to an affine cyclic family requires a separate extraction
and coverage theorem. The two geometric interfaces in this article are
useful complementary components; no unproved adapter between them is
silently assumed.
